% Template naskah akademik FILKOM Universitas Brawijaya.
% Berkas ini hanya merangkai bagian-bagian; SELURUH format ada di gaya/ubtesis.sty.
% Kompilasi: latexmk main.tex   (pdfLaTeX + biber)
\documentclass[12pt,a4paper,oneside]{report}
\usepackage{gaya/ubtesis}
\addbibresource{referensi.bib}

% ------------------------------------------------------------- data sampul
% Ganti seluruh isi di bawah ini. Hanya bagian inilah yang wajib disunting
% sebelum naskah dapat dikompilasi dengan identitas yang benar.
\newcommand\judul{Judul Naskah Ditulis dengan Huruf Kapital di Awal Setiap Kata
  dan Boleh Melampaui Satu Baris}
\newcommand\jenisnaskah{PROPOSAL TESIS}   % atau TESIS, SKRIPSI, PROPOSAL SKRIPSI
\newcommand\penulis{Nama Lengkap Penulis}
\newcommand\nim{000000000000000}
\newcommand\programstudi{PROGRAM STUDI ...}
\newcommand\departemen{DEPARTEMEN ...}
\newcommand\fakultas{FAKULTAS ILMU KOMPUTER}
\newcommand\universitas{UNIVERSITAS BRAWIJAYA}
\newcommand\kota{MALANG}
\newcommand\tahun{20XX}
\newcommand\pembimbing{[NAMA PEMBIMBING]}

\begin{document}
% ------------------------------------------------------------- halaman awal
% Margin atas 3,5 cm dan nomor halaman angka Romawi kecil.
\margininduk
\pagenumbering{roman}
% Sampul [spek D]. Tanpa nomor halaman, tetapi terhitung sebagai halaman i.
% Jarak vertikal disetel terhadap Template Tesis V2.pdf hal-1; lihat _verifikasi/v2/.
% Judul panjang akan menambah baris dan menggeser blok atas ke bawah,
% sama seperti di Word; blok institusi tetap tertambat di bawah.
\begin{titlepage}
\thispagestyle{empty}
\setlength{\parskip}{0pt}
\centering
\vspace*{-0.57cm}  % V2: garis dasar judul 4,04 cm dari tepi atas
{\bfseries\ukuran{16}\judul\par}   % V2: 16 pt tebal, huruf judul (Title Case), bukan kapital
\vspace{0.89cm}
{\bfseries\ukuran{14}\jenisnaskah\par}              % AcademicLevel 14 pt tebal
\vspace{2.55cm}
{\ukuran{14}Disusun oleh:\par
 \penulis\par
 NIM: \nim\par}
\vspace{2.72cm}  % V2: logo 5 cm, tepi atas 14,71 cm
\includegraphics[width=5cm]{logo-ub}\par
\vfill
{\ukuran{14}\programstudi\par
 \departemen\par
 \fakultas\par
 \universitas\par
 \kota\par
 \tahun\par}
\vspace{0.16cm}  % V2: garis dasar tahun 26,54 cm dari tepi atas
\end{titlepage}
\setcounter{page}{2}

\daftarisi
\daftartabel
\daftargambar
\daftarlampiran         % terisi otomatis dari \appendix; hapus bila tanpa lampiran

% ------------------------------------------------------------- isi
% Margin atas kembali 3 cm dan nomor halaman angka Arab.
\marginisi
\clearpage
\pagenumbering{arabic}
% Satu berkas per bab. \chapter menghasilkan "BAB I" lalu judulnya di baris berikut.
\chapter{Pendahuluan}\label{bab:pendahuluan}

\section{Latar Belakang}\label{sec:latar-belakang}

Paragraf pertama setiap bab dan subbab tidak perlu perlakuan khusus; indentasi 0,75~cm
sudah diatur gaya. Rujukan silang ditulis \texttt{Bab~\textbackslash ref\{bab:pendahuluan\}}
sehingga menghasilkan Bab~\ref{bab:pendahuluan}.

Sitasi memakai biblatex gaya \emph{authoryear}: \parencite{contoh2020} untuk sitasi dalam
kurung, dan \textcite{contoh2020} bila nama penulis menjadi bagian kalimat.

Istilah asing dicetak miring, misalnya \textit{deep learning}. Tanda yang masih harus
dilengkapi ditandai agar mudah dicari: \todo{lengkapi angka ini},
\rujukandiperlukan{klaim ini}, dan \konfirmasi{tanyakan ke pembimbing}.

\section{Rumusan Masalah}\label{sec:rumusan-masalah}

\begin{enumerate}[label=RM\arabic*.,leftmargin=1.2cm,labelwidth=1.2cm]
  \item Pertanyaan penelitian pertama.
  \item Pertanyaan penelitian kedua.
\end{enumerate}

\section{Contoh Tabel}\label{sec:contoh-tabel}

Keterangan tabel diletakkan di atas tabel, keterangan gambar di bawah gambar.

\begin{table}[H]
  \centering
  \caption{Keterangan tabel ditulis di atas}
  \label{tab:contoh}
  \begin{tabular}{|l|l|}
    \hline
    \textbf{Kolom pertama} & \textbf{Kolom kedua} \\ \hline
    Isi baris pertama      & Isi baris pertama    \\ \hline
    Isi baris kedua        & Isi baris kedua      \\ \hline
  \end{tabular}
  \sumber{Penulis, \tahun.}
\end{table}

Untuk tabel yang melampaui satu halaman, pakai \texttt{xltabular} dan perintah
\texttt{\textbackslash lanjutan\{n\}} pada baris penanda sambungan.

\section{Contoh Persamaan}\label{sec:contoh-persamaan}

Persamaan bernomor otomatis mengikuti nomor bab.

\begin{equation}\label{eq:contoh}
  y = \sum_{i=1}^{n} w_i x_i + b
\end{equation}

Persamaan~\ref{eq:contoh} dirujuk seperti ini.

% \input{bab/bab2}
% \input{bab/bab3}

% ------------------------------------------------------------- akhir
\daftarreferensi

% Setiap \chapter sesudah \appendix menjadi "LAMPIRAN A ..." dan otomatis
% tercatat di DAFTAR LAMPIRAN. Hapus blok ini bila naskah tanpa lampiran.
\appendix
% Berkas ini di-input sesudah \appendix, sehingga \chapter menjadi "LAMPIRAN A".
\chapter{Judul Lampiran Pertama}\label{lam:a}

Isi lampiran. Tabel dan gambar di sini tetap bernomor mengikuti huruf lampiran,
misalnya Tabel~A.1 dan Gambar~A.1.

\end{document}
